\PassOptionsToPackage{no-math}{fontspec}
\documentclass[UTF8,10pt,twoside,openany,scheme=plain,fontset=fandol]{ctexbook}
\usepackage[paperwidth=155mm,paperheight=235mm,textwidth=332bp,
 top=55bp,bottom=56bp,inner=51bp,headheight=13pt,headsep=10pt,footskip=22pt]{geometry}
\usepackage{amsmath,newtxmath}
\setmainfont{TeX Gyre Termes}
\setsansfont{TeX Gyre Heros}
\setmonofont{TeX Gyre Cursor}
\usepackage{graphicx,tabularx,booktabs,longtable,array,multicol}
\usepackage[table]{xcolor}
\usepackage{enumitem,needspace,fancyhdr,emptypage,caption,fvextra,placeins}
\usepackage[most]{tcolorbox}
\usepackage{etoolbox,xurl}
\usepackage[unicode,hidelinks,bookmarksnumbered]{hyperref}
\usepackage{bookmark}

\hypersetup{pdftitle={深度生成建模（第二版）},pdfauthor={Jakub M. Tomczak}}
\urlstyle{same}
\renewcommand{\contentsname}{目录}
\renewcommand{\figurename}{图}
\renewcommand{\tablename}{表}
\renewcommand{\indexname}{索引}
\ctexset{
  chapter={name={第,章},number=\arabic{chapter},
    format=\raggedright\sffamily\bfseries\fontsize{16}{22}\selectfont,
    beforeskip=0pt,afterskip=106pt},
  section={format=\raggedright\sffamily\bfseries\fontsize{12}{16}\selectfont,
    beforeskip=15pt plus 2pt,afterskip=8pt},
  subsection={format=\raggedright\sffamily\bfseries\fontsize{11}{15}\selectfont,
    beforeskip=13pt plus 2pt,afterskip=6pt},
  subsubsection={format=\raggedright\sffamily\bfseries\fontsize{10.5}{15}\selectfont,
    beforeskip=11pt plus 2pt,afterskip=4pt},
  paragraph={runin=false,format=\raggedright\sffamily\bfseries\fontsize{10.5}{15}\selectfont,
    beforeskip=8pt plus 2pt,afterskip=4pt}
}
\setcounter{secnumdepth}{4}
\setcounter{tocdepth}{4}

\makeatletter
\renewcommand{\normalsize}{%
  \@setfontsize\normalsize{10.5pt}{15pt}%
  \abovedisplayskip 9pt plus 2pt minus 3pt%
  \abovedisplayshortskip 4pt plus 2pt%
  \belowdisplayskip 9pt plus 2pt minus 3pt%
  \belowdisplayshortskip 5pt plus 2pt minus 2pt%
  \let\@listi\@listI}
\renewcommand{\small}{\@setfontsize\small{9.5pt}{13.5pt}}
\renewcommand{\footnotesize}{\@setfontsize\footnotesize{9pt}{12.5pt}}
\newcommand{\dgmcontents}{%
  \chapter*{\contentsname}\markboth{\contentsname}{\contentsname}%
  \pdfbookmark[0]{\contentsname}{dgm-contents}%
  \@starttoc{toc}}
\makeatother
\AtBeginDocument{\normalsize\setlength{\parindent}{2\ccwd}}
\setlength{\parskip}{0pt}
\setlength{\emergencystretch}{1em}
\clubpenalty=7500
\widowpenalty=7500
\raggedbottom
\setlist{leftmargin=2em,topsep=4pt,itemsep=2pt,parsep=0pt}

\pagestyle{fancy}
\fancyhf{}
\fancyhead[LE,RO]{\footnotesize\thepage}
\fancyhead[RE]{\footnotesize\nouppercase{\leftmark}}
\fancyhead[LO]{\footnotesize\nouppercase{\rightmark}}
\renewcommand{\headrulewidth}{0pt}
\renewcommand{\chaptermark}[1]{\markboth{\ifnum\value{chapter}>0 \thechapter\quad\fi #1}{#1}}
\renewcommand{\sectionmark}[1]{\markright{\thesection\quad #1}}
\fancypagestyle{plain}{%
  \fancyhf{}%
  \fancyfoot[LE,RO]{\footnotesize\thepage}%
  \renewcommand{\headrulewidth}{0pt}}

\captionsetup{font=footnotesize,labelfont=bf,labelsep=quad,
  justification=raggedright,singlelinecheck=false,skip=5pt}
\captionsetup[table]{position=top}
\newcolumntype{L}[1]{>{\raggedright\arraybackslash}p{#1}}
\newcolumntype{Y}{>{\raggedright\arraybackslash}X}
\setlength{\tabcolsep}{3.5pt}
\renewcommand{\arraystretch}{1.16}
\definecolor{notegreen}{HTML}{087817}
\newtcolorbox{booknote}[2][]{
  enhanced,breakable,colframe=notegreen,colback=notegreen!5,coltitle=white,
  fonttitle=\sffamily\bfseries,arc=1mm,boxrule=.5pt,title={#2},
  before skip=6pt,after skip=6pt,left=7pt,right=7pt,top=5pt,bottom=5pt,#1}
\newcounter{booklisting}[chapter]
\renewcommand{\thebooklisting}{\thechapter.\arabic{booklisting}}
\newcommand{\bookcode}[3]{%
  \par\medskip
  \VerbatimInput[fontsize=\footnotesize,numbers=left,numbersep=6pt,
    xleftmargin=2em,breaklines,breakanywhere,frame=single,framerule=.2pt,
    rulecolor=\color{black!35}]{#1}%
  \nopagebreak\refstepcounter{booklisting}\label{#3}%
  {\footnotesize\noindent\textbf{代码清单~\thebooklisting}\quad #2\par}%
  \medskip}

\newcommand{\frontchapter}[1]{%
  \chapter*{#1}\markboth{#1}{#1}\addcontentsline{toc}{chapter}{#1}}
\newcommand{\bookimprint}[2][作者]{%
  \insert\footins{\normalfont\footnotesize
    \hsize=\columnwidth\parindent=0pt
    \noindent © #1独家授权 Springer Nature Switzerland AG，2024。%
    \newline J.~M.~Tomczak，《深度生成建模》，%
    \url{#2}。\par}}
\newcommand{\sourcechapter}[1]{%
  \label{chapter:start:#1}%
  \bookimprint{https://doi.org/10.1007/978-3-031-64087-2_#1}}
\newcommand{\citeref}[1]{\hyperref[ref:\thechapter-#1]{[#1]}}
\newcommand{\citerefs}[1]{%
  \begingroup
  \def\refsep{}%
  \def\chaptercite##1{\refsep\hyperref[ref:\thechapter-##1]{##1}\def\refsep{, }}%
  [\forcsvlist{\chaptercite}{#1}]%
  \endgroup}
\newenvironment{chapterreferences}{%
  \section*{参考文献}\markright{参考文献}\addcontentsline{toc}{section}{参考文献}%
  \small\raggedright
  \begin{enumerate}[label=\arabic*.,leftmargin=2.2em,itemsep=3pt]%
}{\end{enumerate}}
\newcommand{\refentry}[1]{\item\label{ref:\thechapter-#1}}
\newcommand{\figref}[1]{图~\ref{#1}}
\newcommand{\tabref}[1]{表~\ref{#1}}
\NewDocumentCommand{\dgmfig}{m m m m o}{%
  \begin{figure}[htbp]
    \centering\includegraphics[width=#2,clip]{figures/#1.pdf}
    \caption{#3}\label{#4}
    \IfValueT{#5}{\diagramlabels{#5}}%
  \end{figure}%
}
\newcommand{\diagramlabels}[1]{\par\smallskip{\raggedright\footnotesize\noindent #1\par}}
